\section{The Standard Face Is a File}
\label{sec:standard-face}
Coded exposure is the hardest of Benjamin's four to hold onto, because both halves of it are real at once and the remedy for either makes the other worse. The worked case is the one named above: face recognition. What follows takes that case apart, because the arithmetic inside it is legible in a way most of this chapter's material is not.

\runin{The machine doing its arithmetic aloud}
Consider what a face pipeline does before it computes anything about a face. DeepFace detects a face and then warps it onto a single generic three-dimensional model — built by averaging a set of scanned human heads — so that every image reaches the network frontal, centered, and in correspondence with the same reference \autocite{taigman2014deepface}. Whatever is learned afterward is learned downstream of that registration. The standard face is not a figure of speech in this arrangement. It is a file, somebody assembled it, and every other face on earth is presented to the system as a departure from it.

Then consider what commercial products in this field put out. Joy Buolamwini and Timnit Gebru put three commercial gender classifiers against a benchmark they built to be balanced across skin phenotype and gender, and the error rates sorted themselves: up to 0.8 percent on lighter-skinned male faces, up to 34.7 percent on darker-skinned female faces \autocite{buolamwini2018gender}. These are different systems. DeepFace is Facebook's and the three audited classifiers are Microsoft's, IBM's and Face++'s; no published result connects the error rates in the second to the registration step in the first. Read as an accuracy result, that is three products wanting better training data. Read as arithmetic, it has the shape of an ordering nobody had to instruct — a face at the center, and everyone else at a measured distance from it. The audit measured two phenotype categories rather than a gradient, so what the pair shows is the ordering, and not a rate at which error climbs with distance from the reference.

The audit's own construction belongs to the finding, and its authors say so. Gender was scored as female or male because the products emit binary sex labels, and the paper states that its evaluation inherits those labels and the reduced view of gender they carry \autocite{buolamwini2018gender}. Skin was scored on the Fitzpatrick scale, a dermatological instrument built to predict sunburn. To measure what the machine was doing, the audit had to speak in the machine's categories. That is not a defect in the work but the cleanest demonstration available that the categories belong to the instrument rather than to the world, and it is why the problem cannot be closed downstream of the instrument.

Balancing the corpus is worth doing and does not reach it. A balanced training set can move the center and can flatten the error rates, which would matter to real people. What it cannot do is remove the reference the pipeline registers against, because that registration is the first operation and something has to occupy the middle of it. The politics gets relocated rather than dissolved: a different face sits at the center, and everyone else is still at a distance, still being measured.

\runin{What is given is a head}
What is given anatomically is a head. A face is something else — a surface organized to be read, held still, framed, lit, and set against other faces of the same kind. Deleuze and Guattari's plateau on the subject makes the strong version of the claim: concrete faces are not found but produced, by what they call an abstract machine of faciality, which overcodes the head and turns it into a surface where meaning is expected to appear \autocite[Plateau~7]{deleuze1987plateaus}. The machine is not itself a face but the apparatus that makes faces, and their account of how it does racism is that it does not proceed by exclusion at all: “Racism operates by the determination of degrees of deviance to the White man's face, which endeavors to integrate nonconforming traits into increasingly eccentric and backward waves, sometimes tolerating them at a given place under given conditions, in a ghetto, or sometimes erasing them from the wall, which never abides alterity” \autocite[178]{deleuze1987plateaus}. Nobody is put outside. Everyone is placed at a measured distance from a center, and the measuring is the operation.

\runin{What the science has not settled}
The applied field rests on a premise its own underlying science has not established: that emotions are the kind of thing a system reads off a signal, the way a thermometer reads temperature.

The case for discrete, evolved emotions rests on Paul Ekman's work with the Fore of Papua New Guinea, who matched short stories to photographs of Western faces much as Westerners did \autocite{ekman1969pancultural}, and on Jaak Panksepp's location of seven primal emotional systems in subcortical structures shared across mammals \autocite{panksepp1998affective}. The rival camp went looking for that circuitry and did not find it: Lisa Feldman Barrett's meta-analysis of imaging studies found no dedicated fear or disgust center, the same regions reappearing across named emotions in shifting combinations \autocite{lindquist2012brain}, and on her account an emotion is a prediction the brain settles on rather than a state it detects. With the forced-choice method dropped, the Himba did not group faces into the Western basic-emotion categories \autocite{gendron2014perceptions}, and Trobriand Islanders read the gasping face Westerners call fear as a threat display \autocite{crivelli2016fear}.

\runin{What neither camp is arguing about}
The field has not resolved that dispute, and the part of it this chapter needs is not in dispute at all. Set the two positions side by side and notice what they leave alone. For Ekman the face displays a discrete internal state; for Barrett it is a poor guide to a state the brain is assembling as it goes. They disagree about what the face discloses and how much survives the trip. Both take the object of the reading as given: there is a face, it was there before anyone looked, and the open question is what can be recovered from it. That premise is the one the applied field actually rests on, and neither camp has been defending it.

If the constructionists are right, the face is a thin signal for feeling. If the faciality argument is right, the readable face is an artifact of the same apparatus that reads it, and a system trained to recognize emotion from faces is trained on that apparatus's output — held still, well lit, correctly registered, and labeled by people who learned the categories from it. A high accuracy score then partly measures how thoroughly the apparatus has already done its work.

\runin{Affective computing, and the four difficulties it does not design away}
The applied half proceeds anyway: systems that read and respond to emotional signal taken from the face, the voice, and the body. Some systems using facial, vocal, or physiological signals may be useful under either account, but usefulness must be demonstrated for a specific task. Neither theory licenses treating those signals as a validated reading of an individual's emotion by default. Facial expression analysis is a harder problem than the face verification deployed at scale, which only decides whether two images are the same person, and accuracy claims that borrow the reputation of the second for the first should be read with that in mind \autocite{taigman2014deepface}. Physiological sensing, which infers from heart rate and skin conductance, is the only one of the three the subject cannot pose for.

Four difficulties travel with the whole field. Emotional expression varies by culture and by person, so a system trained on one population misreads another \autocite{barrett2019emotional,elfenbein2002universality}. Whatever bias is in the training data is in the classifier — which is the smaller of the two claims available here, the larger one being that the norm is not an accident of the sample but a component of the instrument. Consent to emotional inference is not obtained by a term of service. And a system that models emotion well enough to respond to it models it well enough to exploit it.

\runin{What an audit would have to measure}
Audit the registration stage and not only the classifications. Publish the reference model, or enough to reconstruct it, and report error as a function of each face's distance from it, separately for populations represented and unrepresented in training. If error does not rise with distance and stays flat across those populations, the account above is wrong for that system. If it rises, the system's validated use is confined to the region where it was measured, and whether it measures anything but distance is still open. Vendors have no reason to expose a reference model, so this is an audit somebody has to be able to compel.
